% GENERATED FILE. Edit canonical lessons, then run npm run generate:books.
% canonical-source-hash: fnv1a64:a4993f5ee98236d9
% canonical-lessons: LA-C488-bruma, LA-C488-solstitium, LA-C488-canicula, LA-C488-aestivus, LA-C488-caleo, LA-R488-review-487-488, LA-R488-review-final-488

\chapter{Midwinter, Summer solstice, Dog Star, Of summer, To be warm}
\label{ch:la-midwinter-summer-solstice-dog-star-of-summer-to-be-warm}
\hlchaptermodality{488}

\begin{chapteropening}
\textbf{By the end of this chapter:} \emph{I can talk about the turning points of the year: midwinter, the summer solstice, the Dog Days, summer things and warm weather.}

Complete the last lesson of chapter 488: I can talk about the turning points of the year: midwinter, the summer solstice, the Dog Days, summer things and warm weather.
\end{chapteropening}

\section[brūma, brūmae]{brūma, brūmae — midwinter, the winter solstice, the shortest day}

\label{lesson:LA-C488-bruma}



\begin{quote}
Before the new one: say the Latin for to blow through, then the Latin for clear weather.
\end{quote}

\subsection*{You'll want to know: brūma, brūmae}
\textbf{brūma, brūmae} — \textquotedblleft{}midwinter, the winter solstice, the shortest day\textquotedblright{}.

\begin{quote}
Brūmā diēs brevissimus est et nox longissima.
\end{quote}

\begin{quote}
\emph{At midwinter the day is shortest and the night longest.}
\end{quote}

First declension, feminine: \emph{brūma, brūmae f.} Shortened from \emph{brevissima (diēs)}, the shortest day. English \emph{brumal} 'wintry' comes from it.

\begin{cousinweb}
\emph{brūma} comes from \emph{brevima}, an old form of \emph{brevissima} 'shortest', as in \emph{brevissima diēs} 'the shortest day' (EtymDB; Lewis \& Short); it is built on \emph{brevis} 'short'.
\end{cousinweb}

\begin{grammarlens}[title={What you've built}]
One more word for telling, explaining and describing at B1.
\end{grammarlens}

\subsection*{Your turn}
\begin{itemize}
\raggedright
  \item \emph{Say it:} \emph{brūma, brūmae}
  \item \emph{Say it:} \emph{brūma, brūmae}, once more
  \item \emph{Say it:} a sentence of your own with \emph{brūma, brūmae}
  \item \emph{Recall:} say the Latin for to blow through, then the Latin for clear weather, then say \emph{brūma, brūmae} again
\end{itemize}

\begin{tcolorbox}[breakable,colback=teal!4,colframe=teal!35!black,title={Before you move on}]
What does brūma, brūmae mean? (\textquotedblleft{}Midwinter, the winter solstice, the shortest day\textquotedblright{}.) Say it once more.
\end{tcolorbox}
\section[sōlstitium, sōlstitiī]{sōlstitium, sōlstitiī — the summer solstice; midsummer heat}

\label{lesson:LA-C488-solstitium}



\begin{quote}
Before the new one: say the Latin for clear weather, then the Latin for midwinter.
\end{quote}

\subsection*{You'll want to know: sōlstitium, sōlstitiī}
\textbf{sōlstitium, sōlstitiī} — \textquotedblleft{}the summer solstice; midsummer heat\textquotedblright{}.

\begin{quote}
Sōlstitiō diēs longissimus est.
\end{quote}

\begin{quote}
\emph{At the summer solstice the day is longest.}
\end{quote}

Second declension, neuter: \emph{sōlstitium, sōlstitiī n.} For Romans mostly the summer one; midwinter was \emph{brūma}. English \emph{solstice} comes from it.

\begin{cousinweb}
\emph{sōlstitium} is built from \emph{sōl} 'sun' and \emph{sistō} 'to stand still' (EtymDB; Lewis \& Short): the time when the sun seems to stand still. English \emph{solstice} comes from it.
\end{cousinweb}

\begin{grammarlens}[title={What you've built}]
One more word for telling, explaining and describing at B1.
\end{grammarlens}

\subsection*{Your turn}
\begin{itemize}
\raggedright
  \item \emph{Say it:} \emph{sōlstitium, sōlstitiī}
  \item \emph{Say it:} \emph{sōlstitium, sōlstitiī}, once more
  \item \emph{Say it:} a sentence of your own with \emph{sōlstitium, sōlstitiī}
  \item \emph{Recall:} say the Latin for clear weather, then the Latin for midwinter, then say \emph{sōlstitium, sōlstitiī} again
\end{itemize}

\begin{tcolorbox}[breakable,colback=teal!4,colframe=teal!35!black,title={Before you move on}]
What does sōlstitium, sōlstitiī mean? (\textquotedblleft{}The summer solstice; midsummer heat\textquotedblright{}.) Say it once more.
\end{tcolorbox}
\section[canīcula, canīculae]{canīcula, canīculae — the Dog Star; a little dog}

\label{lesson:LA-C488-canicula}



\begin{quote}
Before the new one: say the Latin for midwinter, then the Latin for the summer solstice.
\end{quote}

\subsection*{You'll want to know: canīcula, canīculae}
\textbf{canīcula, canīculae} — \textquotedblleft{}the Dog Star; a little dog\textquotedblright{}.

\begin{quote}
Ubi canīcula ardet, agricolae in umbrā quiēscunt.
\end{quote}

\begin{quote}
\emph{When the Dog Star blazes, the farmers rest in the shade.}
\end{quote}

First declension, feminine: \emph{canīcula, canīculae f.}, a diminutive of \emph{canis} 'dog'. Its rising in high summer brought the hottest days, the dog days.

\begin{cousinweb}
\emph{canīcula} is a diminutive of \emph{canis} 'dog' with \emph{-cula} (EtymDB; Lewis \& Short): a little dog, and the star in the mouth of the Great Dog. English \emph{canicular} comes from it.
\end{cousinweb}

\begin{grammarlens}[title={What you've built}]
One more word for telling, explaining and describing at B1.
\end{grammarlens}

\subsection*{Your turn}
\begin{itemize}
\raggedright
  \item \emph{Say it:} \emph{canīcula, canīculae}
  \item \emph{Say it:} \emph{canīcula, canīculae}, once more
  \item \emph{Say it:} a sentence of your own with \emph{canīcula, canīculae}
  \item \emph{Recall:} say the Latin for midwinter, then the Latin for the summer solstice, then say \emph{canīcula, canīculae} again
\end{itemize}

\begin{tcolorbox}[breakable,colback=teal!4,colframe=teal!35!black,title={Before you move on}]
What does canīcula, canīculae mean? (\textquotedblleft{}The Dog Star; a little dog\textquotedblright{}.) Say it once more.
\end{tcolorbox}
\section[aestīvus, aestīva, aestīvum]{aestīvus, aestīva, aestīvum — of summer, summer}

\label{lesson:LA-C488-aestivus}



\begin{quote}
Before the new one: say the Latin for the summer solstice, then the Latin for the Dog Star.
\end{quote}

\subsection*{You'll want to know: aestīvus, aestīva, aestīvum}
\textbf{aestīvus, aestīva, aestīvum} — \textquotedblleft{}of summer, summer\textquotedblright{}.

\begin{quote}
Dīves vīllam aestīvam in montibus habet.
\end{quote}

\begin{quote}
\emph{The rich man has a summer house in the mountains.}
\end{quote}

An adjective of three endings, from \emph{aestās} 'summer'. \emph{Aestīva} (n. pl.), a summer camp. English \emph{aestivate} comes from the same root.

\begin{cousinweb}
\emph{aestīvus} is built on \emph{aestās} 'summer' with the suffix \emph{-īvus} (EtymDB; Lewis \& Short), and \emph{aestās} goes back to PIE \emph{h\textsubscript{2}eyd\textsuperscript{h}-} 'to burn'.
\end{cousinweb}

\begin{grammarlens}[title={What you've built}]
One more word for telling, explaining and describing at B1.
\end{grammarlens}

\subsection*{Your turn}
\begin{itemize}
\raggedright
  \item \emph{Say it:} \emph{aestīvus, aestīva, aestīvum}
  \item \emph{Say it:} \emph{aestīvus, aestīva, aestīvum}, once more
  \item \emph{Say it:} a sentence of your own with \emph{aestīvus, aestīva, aestīvum}
  \item \emph{Recall:} say the Latin for the summer solstice, then the Latin for the Dog Star, then say \emph{aestīvus, aestīva, aestīvum} again
\end{itemize}

\begin{tcolorbox}[breakable,colback=teal!4,colframe=teal!35!black,title={Before you move on}]
What does aestīvus, aestīva, aestīvum mean? (\textquotedblleft{}Of summer, summer\textquotedblright{}.) Say it once more.
\end{tcolorbox}
\section[caleō, calēre]{caleō, calēre — to be warm, to be hot}

\label{lesson:LA-C488-caleo}



\begin{quote}
Before the new one: say the Latin for the Dog Star, then the Latin for of summer.
\end{quote}

\subsection*{You'll want to know: caleō, calēre}
\textbf{caleō, calēre} — \textquotedblleft{}to be warm, to be hot\textquotedblright{}.

\begin{quote}
Aestāte saxa sub sōle calent.
\end{quote}

\begin{quote}
\emph{In summer the stones are hot under the sun.}
\end{quote}

\emph{caleō, calēre, caluī}, second conjugation. \emph{Calidus} 'hot' and \emph{calor} 'heat' come from it; English \emph{calorie} comes from \emph{calor}.

\begin{cousinweb}
\emph{caleō} is inherited from Proto-Italic \emph{kalēō} 'to be hot', from PIE \emph{\'{k}elh\textsubscript{1}-} 'to heat' (EtymDB). \emph{calidus} 'hot' and \emph{calor} 'heat', the source of English \emph{calorie}, are built on it (Lewis \& Short).
\end{cousinweb}

\begin{grammarlens}[title={What you've built}]
One more word for telling, explaining and describing at B1.
\end{grammarlens}

\subsection*{Your turn}
\begin{itemize}
\raggedright
  \item \emph{Say it:} \emph{caleō, calēre}
  \item \emph{Say it:} \emph{caleō, calēre}, once more
  \item \emph{Say it:} a sentence of your own with \emph{caleō, calēre}
  \item \emph{Recall:} say the Latin for the Dog Star, then the Latin for of summer, then say \emph{caleō, calēre} again
\end{itemize}

\begin{tcolorbox}[breakable,colback=teal!4,colframe=teal!35!black,title={Before you move on}]
What does caleō, calēre mean? (\textquotedblleft{}To be warm, to be hot\textquotedblright{}.) Say it once more.
\end{tcolorbox}
\section[(dialogue)]{Farm, mountains and weather}

\label{lesson:LA-R488-review-487-488}



\begin{quote}
Say the Latin for the north wind and the south wind.
\end{quote}

\subsection*{Your turn}
Say these aloud:

\begin{itemize}
\raggedright
  \item the north wind, the south wind, the west wind, to blow through, clear weather
  \item midwinter, the summer solstice, the Dog Star, of summer, to be warm
\end{itemize}

\begin{tcolorbox}[breakable,colback=teal!4,colframe=teal!35!black,title={Before you move on}]
The north wind? (\textbf{aquilō, aquilōnis}.) The south wind? (\textbf{auster, austrī}.) The west wind? (\textbf{favōnius, favōniī}.) To blow through? (\textbf{perflō, perflāre}.) Clear weather? (\textbf{serēnitās, serēnitātis}.) Midwinter? (\textbf{brūma, brūmae}.) The summer solstice? (\textbf{sōlstitium, sōlstitiī}.) The Dog Star? (\textbf{canīcula, canīculae}.) Of summer? (\textbf{aestīvus, aestīva, aestīvum}.) To be warm? (\textbf{caleō, calēre}.)
\end{tcolorbox}
\section[(dialogue)]{Midwinter, Summer solstice, Dog Star, Of summer, To be warm — once more}

\label{lesson:LA-R488-review-final-488}



\begin{quote}
Say the Latin for midwinter and the summer solstice.
\end{quote}

\subsection*{Your turn}
\begin{itemize}
\raggedright
  \item \emph{Say it:} midwinter, the summer solstice, the Dog Star, of summer, to be warm
\end{itemize}

\begin{tcolorbox}[breakable,colback=teal!4,colframe=teal!35!black,title={Before you move on}]
Midwinter? (\textbf{brūma, brūmae}.) The summer solstice? (\textbf{sōlstitium, sōlstitiī}.) The Dog Star? (\textbf{canīcula, canīculae}.) Of summer? (\textbf{aestīvus, aestīva, aestīvum}.) To be warm? (\textbf{caleō, calēre}.)
\end{tcolorbox}
